\scanpage{68}
\mspage{55}
% archive sheet number: 67

qu'on soit \inserted{pour $M,W,C,I$} dans l'une des conditions a), b) du th. 3.1. Alors $u_!$ induit une flèche

\[
 u_{!\,\mathrm{cof}}:M(I^{\unsure{\circ|0}})_C\mathbin{-} M(I'^\circ)_C\subset M(I'^\circ),
\]

qui transforme flèches de $W(I^\circ)$ en flèches de $W(I'^\circ)$ (i.e compatible avec les localiseurs), Il induit donc un foncteur sur les localisées, d'où \deleted{par} \inserted{un} \deleted{$I'$ satisfaisant \unsure{égalt}} foncteur \deleted{foncteur}

\[
 u_!^W:\mathrm{Ho}_W(I^\circ)\longrightarrow\mathrm{Ho}_W(I'^\circ).\tag*{(3.3.1)}
\]

% slanted
\margin{NB Cette flèche est définie, $=$ sans supposer $I'$ ordonné, car $u_{!\,\mathrm{cof}}$ transforme \unsure{une flèche|flèches} de $W(I^\circ)$ en $W(I'^\circ)$. Cf. \unsure{aussi|mém.|même} p.~\unsure{86|88|56}}

\uline{Dém}. Par le th. 3.1, on sait que $(M(I^\circ),W(I^\circ),C\deleted{\unsure{(I^\circ)}}_{I^\circ})$ est une structure de cat. de modèles à cofibration. On sait par ailleurs que $TC_{I^\circ}=C_{I^\circ}\cap W(I^\circ)$, par Th. 2.1. $3^\circ$ (p.~27). Donc par 3.\unsure{2|3}, $u_!$ applique $C_{I^\circ}\cap W(I^\circ)$ dans $TC_{I'^\circ}\subset C_{I'^\circ}\cap\unsure{W(I'^\circ)}$, et a fortiori, en des flèches $\in W(I'^\circ)$.

Par ailleurs, par la théorie des cat à cofibrations, on trouve que toute flèche \inserted{$f$} \deleted{de $W(I^\circ)$ qui est dans} entre objets cofibrants est composée $f=pi$, où $i\in C_{I^\circ}$, et \deleted{deux flèches qui \unsure{est}} où $p$ est $\in W(I^\circ)$ \inserted{et de plus} \uline{inverse à gauche} d'une flèche de $C_{I^\circ}\cap W(I^\circ)$, \unsure{\deleted{de la}|donc} d'une flèche de $(TC)_{I^\circ}$. Alors
